238 \hfill 13 Line and Surface Integrals
\vspace{12pt}
d) Verify that if $f : D \to \mathbb{R}$ and $\mathbf{u} : D \to \mathbb{R}^n$ are smooth functions in a regular
domain $D \subset \mathbb{R}^n$ and $\mathbf{u}$ has no critical points in $D$, then
\[ \int_{D} f \, d\mathbf{v} = \int_{\mathbf{u}^{-1}(\mathbf{t})} \frac{f}{|\nabla \mathbf{u}|} \, d\sigma. \]
e) Let $V_f(t)$ be the measure (volume) of the set $\{x \in D \mid f(x) > t\}$, and let the
function $f$ be nonnegative and bounded in the domain $D$.
Show that $\int_D f \, dV = \int_0^\infty \text{meas} \{V_f(t)\} \, dt = \int_0^\infty V_f(t) \, dt$.
f) Let $\varphi \in C^{(1)}(\mathbb{R}, \mathbb{R}^+)$ and $\varphi(0) = 0$, while $f \in C^{(1)}(D, \mathbb{R})$ and $V_{|f|}(t)$
is the measure of the set $\{x \in D \mid |f(x)| > t\}$. Verify that $\int_D \varphi \circ f \, dV =
\int_0^\infty \varphi'(t) V_{|f|}(t) \, dt$.
\vspace{12pt}
\noindent \textbf{13.3 The Fundamental Integral Formulas of Analysis}
\vspace{12pt}
The most important formula of analysis is the Newton–Leibniz formula (funda-
mental theorem of calculus). In the present section we shall obtain the formulas of
Green, Gauss-Ostrogradskii, and Stokes, which on the one hand are an extension of
the Newton-Leibniz formula, and on the other hand, taken together, constitute the
most-used part of the machinery of integral calculus.
In the first three subsections of this section, without striving for generality in our
statements, we shall obtain the three classical integral formulas of analysis using
visualizable material. They will be reduced to one general Stokes formula in the
fourth subsection, which can be read formally independently of the others.
\vspace{12pt}
\noindent \textbf{13.3.1 Green's Theorem}
\vspace{12pt}
Green's theorem is the following.
\vspace{12pt}
\noindent \textbf{Proposition 1} Let $\mathbb{R}^2$ be the plane with a fixed coordinate grid $x, y$, and let $\bar{D}$ be
a compact domain in this plane bounded by piecewise-smooth curves. Let $P$ and $Q$
be smooth functions in the closed domain $\bar{D}$. Then the following relation holds:
\[ \iint_{\bar{D}} \left( \frac{\partial Q}{\partial x} - \frac{\partial P}{\partial y} \right) dx \, dy = \oint_{\partial \bar{D}} P dx + Q dy, \tag{13.25} \]
in which the right-hand side contains the integral over the boundary $\partial \bar{D}$ of the
domain $\bar{D}$ oriented consistently with the orientation of the domain $\bar{D}$ itself.
\vspace{12pt}
\noindent $^7$G. Green (1793–1841) – British mathematician and mathematical physicist. Newton's grave in
Westminster Abbey is framed by five smaller gravestones with brilliant names: Faraday, Thomson
(Lord Kelvin), Green, Maxwell, and Dirac.